\subsection{有限维自由模自同构的刻画}

定理1. 设M是有限维自由A-模，u是M的一个自同态。以下条件等价：

\begin{enumerate}
\def\labelenumi{(\alph{enumi})}
\item
  u 是双射；
\item
  \(u\) 是右可逆的（II，§1，第 9 号，命题15的推论1）；
\item
  u是左可逆的（II，§1，第9号，命题15的推论2）；
\item
  u 是满射；
\item
  det u 在 A 中可逆。
\end{enumerate}

设 \((e_i)_{1\leqslant i\leqslant n}\) 是 M 的一组基。如果 \(x_{i} = u(e_{i})\) 对于 \(1\leqslant i\leqslant n\) ，则

\[
x _ {1} \wedge x _ {2} \wedge \dots \wedge x _ {n} = (\det u) e _ {1} \wedge e _ {2} \wedge \dots \wedge e _ {n}.
\]

由 §7 第 9 号定理 2 可知， \(x_{i}\) 构成 M 的一组基的必要且充分条件是 \(\det u\) 为 A 的可逆元；这就证明了 (a) 与 (e) 的等价性。注意 (a) 显然蕴含条件 (b)、(c)、(d) 中的每一个；剩下要证的是 (b)、(c)、(d) 中的每一个都蕴含 (e)。现在，若存在 M 的自同态 \(v\) 使得 \(v \circ u = 1_{\mathbf{M}}\) 或 \(u \circ v = 1_{\mathbf{M}}\) ，则 \((\det v)(\det u) = 1\) ，因此 \(\det u\) 在 A 中可逆。若 \(u\) 是满射，则 \(\bigwedge^{n}(u)\) 也是满射（ \(\S 7\) ，第 2 号，命题 3），换言之，A 中比率为 det \(u\) 的位似是满射，这立即蕴含 det \(u\) 可逆。

命题3. 设M是有限维自由A-模。对于M的每个自同态u，下列条件是等价的：

\begin{enumerate}
\def\labelenumi{(\alph{enumi})}
\setcounter{enumi}{5}
\item
  u是单射；
\item
  det u在A中不是零因子。
\end{enumerate}

沿用定理 1 证明中的记号，u 是单射当且仅当 \(x_{i}\) 线性无关。根据 §7 第 9 号命题 12，其必要且充分条件是关系 \(\lambda x_{1} \wedge x_{2} \wedge \cdots \wedge x_{n} = 0\) （其中 \(\lambda \in A\) ）蕴含 \(\lambda = 0\) 。但这等价于 \(\lambda(\det u) = 0\) ，因为 \(e_{1} \wedge \cdots \wedge e_{n}\) 是 \(\bigwedge^{n}(M)\) 的一个基；由此得出命题。

注记. 当A是域时，定理2的条件(e)等价于命题3的条件(g)，因为它们都表示det \(u \neq 0\) 。因此在这种情况下，定理1和命题3的所有条件都是等价的（参见II，§7，第4号，命题9的推论）。

